% Spectre infrarouge de l'éthanol (L'essentiel, chapitre 11)
\begin{tikzpicture}
\begin{axis}[otIR,width=108mm,height=46mm,title={éthanol \ce{CH3-CH2-OH}}]
\addplot[otSpectre] {100 \bandeIR{3350}{170}{62} \bandeIR{2950}{55}{48} \bandeIR{1420}{55}{22} \bandeIR{1080}{45}{55} \bandeIR{880}{30}{18}};
\node[pcLegende,text=otGreenDark,anchor=south] at (axis cs:3350,20) {O--H : large};
\node[pcLegende,text=otNavy,anchor=south] at (axis cs:2820,42) {C--H};
\node[pcLegende,text=otNavy,anchor=south] at (axis cs:1080,32) {C--O};
\end{axis}
\end{tikzpicture}
\par\smallskip\legende{La bande forte et large vers 3350 cm$^{-1}$ signe la liaison O–H d'un alcool ; les bandes C–H apparaissent juste en dessous de 3000 cm$^{-1}$.}
